\documentclass[10pt]{article}

\usepackage[margin=0.85in]{geometry}
\usepackage{amsmath}
\usepackage{booktabs}
\usepackage{microtype}

\title{A Small Demonstration of Agent-Assisted Paper Revision}
\author{Demo Research Group}
\date{}

\begin{document}
\maketitle

\begin{abstract}
This fictional paper demonstrates a review loop in which a reader comments on
the rendered PDF and a coding agent revises the underlying LaTeX source. The
example is intentionally small, self-contained, and unrelated to any real
submission.
\end{abstract}

\section{Introduction}
Technical writing improves through repeated comparison between author intent
and reader interpretation. Source-only review misses layout problems, while
PDF-only review makes precise edits cumbersome. A useful workflow therefore
keeps the rendered document and its source connected without turning the PDF
viewer into a second editor.

The demonstration focuses on three actions: selecting rendered text, discussing
the selection with an agent, and checking the rebuilt PDF after a source edit.

\section{Method}
Let $q_i$ denote the clarity score of paragraph $i$. We summarize a revision by
the average change
\begin{equation}
  \Delta Q = \frac{1}{N}\sum_{i=1}^{N}\left(q_i^{\mathrm{new}}-q_i^{\mathrm{old}}\right).
\end{equation}
The quantity is illustrative; the demo does not claim that writing quality can
be reduced to one metric. Its purpose is to give the PDF a realistic equation,
table, and multi-section structure for interface screenshots.

\begin{table}[h]
  \centering
  \caption{Fictional revision outcomes used in the demonstration.}
  \begin{tabular}{lcc}
    \toprule
    Revision & Comments addressed & Build status \\
    \midrule
    Initial pass & 2 & Success \\
    Follow-up pass & 1 & Success \\
    \bottomrule
  \end{tabular}
\end{table}

\section{Discussion}
The reader remains in control of the paper. The agent receives the selected
quote and the written comment, inspects nearby source, proposes or applies an
edit, compiles the document, and records what changed. Additional comments can
challenge the edit or ask for another revision.

\section{Conclusion}
The demo provides a neutral document for trying tex-mcp-web without exposing a
real manuscript. It also makes the README screenshots reproducible.

\end{document}
